% Part: methods
% Chapter: induction
% Section: induction-on-N

\documentclass[../../../include/open-logic-section]{subfiles}

\begin{document}

\olfileid[pt]{mth}{ind}{inN}

\olsection{Indução em~$\Nat$}

Em sua forma mais simples, a indução é uma técnica usada para provar resultados
para todos os números naturais. Ela usa o fato de que, começando em $0$
e somando~$1$ repetidamente, eventualmente alcançamos todo número natural. Assim,
para provar que algo é verdadeiro para todo número, podemos (1)~estabelecer
que é verdadeiro para $0$ e (2)~mostrar que, sempre que é verdadeiro para um
número~$n$, também é verdadeiro para o número seguinte~$n+1$.  Se
abreviarmos ``o número~$n$ tem a propriedade $P$'' por $P(n)$ (e ``o número~$k$
tem a propriedade $P$'' por $P(k)$, etc.), então uma prova por indução de que
$P(n)$ para todo $n \in \Nat$ consiste em:
\begin{enumerate}
\item uma prova de $P(0)$, e
\item uma prova de que, para qualquer~$k$, se $P(k)$ então $P(k+1)$.
\end{enumerate}
Para deixar isso bem claro, suponha que tenhamos tanto (1) quanto~(2).  Então
(1) nos diz que $P(0)$ é verdadeiro.  Se também tivermos~(2), sabemos em
particular que, se $P(0)$, então $P(0+1)$, isto é, $P(1)$. Isso decorre
da afirmação geral ``para qualquer~$k$, se $P(k)$ então $P(k+1)$'' ao
pôr~$0$ no lugar de~$k$. Então, por modus ponens, temos que~$P(1)$. De (2)
novamente, agora tomando $1$ por~$n$, temos: se $P(1)$ então~$P(2)$. Como
acabamos de estabelecer~$P(1)$, por modus ponens, temos~$P(2)$. E assim
por diante.  Para qualquer número~$n$, depois de fazer isso $n$~vezes, eventualmente
chegamos a~$P(n)$.  Então (1) e~(2) juntos estabelecem~$P(n)$ para qualquer $n
\in \Nat$.

Vejamos um exemplo.  Suponha que queiramos descobrir quantas
somas diferentes podemos lançar com $n$ dados.  Embora possa parecer
bobo, vamos começar com $0$ dados.  Se você não tem dados, há apenas uma
soma possível que você pode ``lançar'': nenhum ponto, que soma~$0$. Assim,
o número de lançamentos possíveis diferentes é~$1$. Se você tem apenas um
dado, isto é, $n=1$, há seis valores possíveis, de $1$ a~$6$. Com
dois dados, podemos lançar qualquer soma de $2$ a $12$, ou seja, $11$
possibilidades.  Com três dados, podemos lançar qualquer número de $3$ a
$18$, isto é, $16$ possibilidades diferentes.  $1$, $6$, $11$, $16$: parece
um padrão: talvez a resposta seja $5n+1$?  É claro que $5n+1$ é o
máximo possível, pois há apenas $5n+1$ números entre $n$,
o menor valor que você pode lançar com $n$ dados (todos $1$) e $6n$, o
maior que você pode lançar (todos $6$).

\begin{thm}
  Com $n$ dados pode-se lançar todos os $5n+1$ valores possíveis entre $n$
  e $6n$.
\end{thm}

\begin{proof}
Seja $P(n)$ a afirmação: ``É possível lançar qualquer número entre
$n$ e~$6n$ usando $n$~dados.''  Para usar a indução, provamos:
\begin{enumerate}
\item A \emph{base da indução} $P(1)$, isto é, com apenas um dado,
  você pode lançar qualquer número entre $1$ e $6$.
\item O \emph{passo da indução}, para todo $k$, se $P(k)$ então~$P(k+1)$.
\end{enumerate}

(1) é provado inspecionando um dado de $6$ lados. Ele tem todos os 6 lados, e
todo número entre $1$ e~$6$ aparece em um dos lados. Então é
possível lançar qualquer número entre $1$ e~$6$ usando um único
dado.

Para provar (2), supomos o antecedente do condicional, isto é,
$P(k)$. Essa suposição é chamada de \emph{hipótese indutiva}.  Nós a
usamos para provar~$P(k+1)$. A parte difícil é encontrar um modo de pensar
sobre os valores possíveis de um lançamento de $k+1$ dados em termos dos
valores possíveis de lançamentos de $k$~dados, somados aos lançamentos do dado
extra, o $(k+1)$-ésimo---é isso, porém, que temos que fazer, se quisermos usar
a hipótese indutiva.

A hipótese indutiva diz que podemos obter qualquer número entre $k$
e~$6k$ usando $k$~dados.  Se lançamos um~$1$ com nosso $(k+1)$-ésimo dado,
isso acrescenta $1$ ao total. Então podemos lançar qualquer valor entre $k+1$
e $6k+1$ lançando $k$~dados e depois rolando um~$1$ com o
$(k+1)$-ésimo dado.  O que falta?  Os valores de $6k+2$ a $6k+6$.  Podemos
obtê-los rolando $k$ dados em $6$ e depois um número entre $2$ e
$6$ com nosso $(k+1)$-ésimo dado. Juntos, isso significa que, com $k+1$ dados,
podemos lançar qualquer um dos números entre $k+1$ e $6(k+1)$, isto é,
provamos~$P(k+1)$ usando a suposição~$P(k)$, a hipótese
indutiva.
\end{proof}

Muito frequentemente usamos a indução quando queremos provar algo sobre uma
série de objetos (números, conjuntos, etc.) que é, ela própria, definida
``indutivamente,'' isto é, definindo o $(n+1)$-ésimo objeto em termos do
$n$-ésimo.  Por exemplo, podemos definir a soma~$s_n$ dos números naturais
até~$n$ por
\begin{align*}
  s_0 & = 0\\
  s_{n+1} & = s_n + (n+1)
\end{align*}
Esta definição dá:
\begin{align*}
  s_0 & = 0,\\
  s_1 & = s_0 + 1 && = 1,\\
  s_2 & = s_1 + 2 && = 1 + 2 = 3\\
  s_3 & = s_2 + 3 && = 1 + 2 + 3 = 6, \text{ etc.}
\end{align*}
Agora podemos provar, por indução, que $s_n = n(n+1)/2$.

\begin{prop}
  $s_n = n(n+1)/2$.
\end{prop}

\begin{proof}
  Temos que provar (1) que $s_0 = 0\cdot(0 + 1)/2$ e (2) se $s_k =
  k(k+1)/2$ então $s_{k+1} = (k+1)(k+2)/2$.  (1) é óbvio. Para provar
  (2), supomos a hipótese indutiva: $s_k = k(k+1)/2$. Usando-a,
  temos que mostrar que $s_{k+1} = (k+1)(k+2)/2$.

  O que é $s_{k+1}$?  Pela definição, $s_{k+1} = s_k + (k+1)$.  Pela
  hipótese indutiva, $s_k = k(k+1)/2$. Podemos substituir isso na
  equação anterior, e então só precisamos de um pouco de aritmética de
  frações:
  \begin{align*}
    s_{k+1} & = \frac{k(k+1)}{2} + (k+1) = {}\\
    & = \frac{k(k+1)}{2} + \frac{2(k+1)}{2} = {}\\
    & = \frac{k(k+1) + 2(k+1)}{2} = {}\\
    & = \frac{(k+2)(k+1)}{2}.
  \end{align*}
\end{proof}

A lição importante aqui é que, se você está provando algo sobre
alguma sequência definida indutivamente $a_n$, a indução é o caminho óbvio
a seguir. E, mesmo quando não é (como no caso das possibilidades de
lançamentos de dados), você pode usar a indução se conseguir, de alguma forma, relacionar o caso
para~$k+1$ ao caso para~$k$.

\end{document}
